\documentclass{article}
\usepackage[T2A]{fontenc}
\usepackage[utf8]{inputenc}
\usepackage[russian]{babel}
\usepackage[margin=2cm]{geometry}
\usepackage{graphicx}
\usepackage{wrapfig}
\usepackage{parskip}
\usepackage{enumitem}
\usepackage{bm}

\setlength{\parindent}{0pt}
\setlength{\parskip}{\baselineskip}

\setcounter{section}{-1}




\begin{document}

\section{Общие сведения о проверке гипотез}

Статистическая гипотеза --- это любое предположение о виде распределения или его параметрах. Например, гипотеза может утверждать, что среднее значение признака равно заданному числу, или что распределение подчиняется нормальному закону.

Если бы нам была известна функция распределения, то любую ее характеристику можно было бы точно посчитать. Но на практике нам известна лишь выборка --- численная реализация случайной величины, которая не может в точности охарактеризовать всю генеральную совокупность и функцию распределения из‑за случайности. Из-за случайности невозможно дать точный ответ об истинности гипотезы, все ответы даются с какой-то вероятностью.

Чтобы перейти от абстрактных рассуждений к конкретной процедуре, используют статистический критерий. Критерий --- это формальное правило, которое включает в себя статистику (специальную функцию от выборки, значение которой является числом), и ее теоретическое распределение при условии, что нулевая гипотеза верна (то есть для выборок, которые удовлетворяют гипотезе и соответсвенно подобранному критерию). Статистика устроена так, что её значение численно отражает меру расхождения между наблюдаемыми данными и тем, что предсказывает $H_0$. Любое наблюдаемое отклонение можно оценить на графике её распределения: если значение статистики попадает в хвосты — области, вероятность попадания в которые при верной $H_0$ крайне мала, — это служит свидетельством того, что данные не согласуются с нулевой гипотезой. Таким образом, критическая область критерия напрямую отражает, насколько сильно фактические данные противоречат $H_0$.

\begin{wrapfigure}{r}{0.25\textwidth}
    \centering
    \includegraphics[width=\linewidth]{image.png}
\end{wrapfigure}

На графике изображена плотность распределения критерия. Чем больше вероятность истинности $H_0$, тем ближе к пику (по горизонтальной оси) будет значение статистики критерия. Красным отмечены области, где вероятность совершить ошибку 1-го рода, то есть отклонить $H_0$, когда она верна, мала. Размер областей зависит от заранее выбранного уровня значимости $\alpha$, вид областей зависит от формулировки альтернативной гипотезы (на рисунке проверяется отклонения от значения в любую из сторон, поэтому критические области симметричны). Уровень значимости $\alpha$ выбирается, как правило, до подсчета статистики.

Чем дальше значение статистики находится от пика распределения, тем менее вероятно, что верна $H_0$, однако она все равно может быть верной, если нам "не повезло" с выборкой. $p$-значение оценивает вероятность, при условии, что $H_0$ верна, получить еще более экстремальный результат, графически это площадь хвостов.

Алгоритм проверки гипотез:
\begin{enumerate}[topsep=-\baselineskip, partopsep=2pt, itemsep=2pt, parsep=0pt]
    \item Сформулировать нулевую и альтернативную гипотезы $H_0$ и $H_1$ (в лабораторной даны)
    \item Выбрать уровень значимости $\alpha$ (в лабораторной дан)
    \item Выбрать критерий (в лабораторной дан)
    \item Посчитать статистику критерия
    \item По распределению критерия найти критическую область или $p-value$ для выбранного уровня значимости (в лабораторной нужно выполнить два варианта)
    \item Принять или отклонить гипотезу
\end{enumerate}

% ==============================================================================

\section{Дать определение вероятности ошибки 1-го рода}

Ошибка 1-го рода совершается, если критерий отверает нулевую гипотезу, когда она верна, то есть мы делаем вывод о том, что различия или эффект есть, хотя на самом деле их нет.

% ==============================================================================

\section{Дать определение вероятности ошибки 2-го рода}

Ошибка 1-го рода совершается, если критерий принимает нелевую гипотезу, когда верная альтернативная, то есть мы пропускаем реального эффекта.

% ==============================================================================

\section{Дать определение размера критерия}

Для простых гипотез $H_0$ можно считать, что размер критерия это уровень значимости $\alpha$. Уровень значимости выбирается заранее (как правило 0.05) и определяет максимальную вероятность, с которой можно отвергнуть $H_0$, когда она на самом деле будет верна. В общем случае, размер критерия --- максимальная вероятность ошибка 1-го рода, которую допускает критерий.

% ==============================================================================

\section{Что означает, что критерий является критерием уровня $\bm\alpha$?}

Это означает, что вероятность соверщить ошибку 1-го рода (отвергнуть $H_0$, когда она верна) не превышает заданного значения уровня значимости $\alpha$.

% ==============================================================================

\section{Что такое $\bm p$-значение (критический уровень значимости)? Чем он отличается от уровня значимости?}

$p$-значение --- это вероятность получить, при верной $H_0$, такой же или еще более крайнее значение ститистики критерия. Пусть имеется конкретная выборка, по которой "видно", что $H_0$ не верна, $p$-значение показывает вероятность получить выборку "еще хуже", если верна $H_0$. Уровень значимости $\alpha$ устанавливается исследователем заранее, далее проверяется попадает ли выборка в соответствующую критическую область, $p$-значение считается для конкретной выборки.

\end{document}